\section{Sugestões para trabalhos futuros}
Recomenda-se caracterizar quimicamente a SPE de \textit{T. chuii} e verificar se os compostos selecionados para prospecção estão presentes no material obtido nas condições de cultivo utilizadas. Essa análise pode identificar os açúcares, derivados sulfatados e fenólicos e quantificar os componentes detectados.

Os compostos selecionados pela comparação de contatos podem ser avaliados em ensaios bioquímicos específicos para cada receptor, incluindo curvas de concentração e controles adequados. Os ensaios permitirão medir alterações na atividade enzimática e comparar os resultados com os perfis previstos por docking.

Como etapa seguinte, a SPE e os compostos selecionados podem ser estudados em modelos celulares da pele, avaliando viabilidade celular, marcadores de estresse oxidativo, componentes da matriz extracelular, hidratação e melanogênese. A parte computacional pode ser ampliada com a validação dos protocolos por redocking, a investigação da flexibilidade dos receptores e o refinamento do tratamento dos centros metálicos.
